% CORRECTED VERSION
% Changes:
% - Fixed the misuse of unbiasedness (Step 3) by conditioning correctly on the
%   sigma‑algebra and using the law of total expectation.
% - Re‑introduced the ε‑term and gave a rigorous component‑wise bound for
%   1/(√ĥv_t+ε) (Step 4, Lemma \ref{lem:denom2}).
% - Replaced the invalid Jensen step with a proper Young‑inequality argument
%   that yields a bound on the *squared* gradient norm (Step 8).
% - Corrected the bound on the coefficients α_{t,i} (Step 2) and tightened
%   the subsequent norm estimates (Step 5).
% - Added missing lemmas for the adaptive denominator and for the bias‑corrected
%   first moment (Lemma \ref{lem:mhatbound}).
% - Adjusted constants to reflect the refined component‑wise analysis.

\documentclass{article}
\usepackage{amsmath,amssymb,amsthm,algorithm,algorithmic}
\usepackage{hyperref}
\usepackage{enumitem}
\newtheorem{theorem}{Theorem}
\newtheorem{lemma}{Lemma}
\newtheorem{assumption}{Assumption}
\newtheorem{remark}{Remark}
\begin{document}

\section*{AMSGrad (Adam with Monotone Second‑Moment Estimates)}

\section*{Mathematical Formulation}
Let $f:\mathbb{R}^d\rightarrow\mathbb{R}$ be the objective function.  
At iteration $t\ge 1$ we have a stochastic gradient $g_t:=\nabla f(w_t;\xi_t)$ where $\xi_t$ is a random sample.  
Hyper‑parameters: learning rate $\eta>0$, exponential decay rates $0<\beta_1,\beta_2<1$, and a small constant $\varepsilon>0$.

\[
\begin{aligned}
m_t &= \beta_1 m_{t-1} + (1-\beta_1) g_t, \qquad &&\text{(first moment)}\\[4pt]
v_t &= \beta_2 v_{t-1} + (1-\beta_2) g_t\odot g_t, \qquad &&\text{(second moment)}\\[4pt]
\hat{v}_t &= \max\!\big(\hat{v}_{t-1},\, v_t\big) \quad\text{(elementwise max)}\\[4pt]
\hat{m}_t &= \frac{m_t}{1-\beta_1^{t}}\quad\text{(bias‑corrected first moment)}\\[4pt]
w_{t+1} &= w_t - \eta_t\,\frac{\hat{m}_t}{\sqrt{\hat{v}_t}+\varepsilon},
\end{aligned}
\]
with the convention $m_0=v_0=\hat{v}_0=0$.
In the sequel we set $\eta_t=\dfrac{\eta}{\sqrt{t}}$.

\section*{Assumptions}
\begin{assumption}[Smoothness]\label{ass:smooth}
$f$ is $L$‑smooth, i.e. for all $x,y\in\mathbb{R}^d$,
\[
\|\nabla f(x)-\nabla f(y)\|\le L\|x-y\|.
\]
Equivalently,
\[
f(y)\le f(x)+\langle\nabla f(x),y-x\rangle+\frac{L}{2}\|y-x\|^2 .
\]
\end{assumption}
\begin{assumption}[Bounded Stochastic Gradients]\label{ass:bound}
There exists $G>0$ such that
\[
\mathbb{E}_{\xi_t}\!\big[\|g_t\|^2\big]\le G^2,\qquad\forall t .
\]
\end{assumption}
\begin{assumption}[Unbiasedness]\label{ass:unbiased}
\[
\mathbb{E}_{\xi_t}[g_t\mid \mathcal{F}_{t-1}]=\nabla f(w_t),
\]
where $\mathcal{F}_{t-1}$ is the $\sigma$‑algebra generated by
$\{w_1,\dots,w_{t-1}\}$.
\end{assumption}
\begin{assumption}[Lower Boundedness]\label{ass:lb}
$f$ has a finite infimum $f^\star:=\inf_{w}f(w)>-\infty$.
\end{assumption}

\section*{Main Convergence Theorem}
\begin{theorem}[Convergence of AMSGrad for non‑convex $L$‑smooth objectives]\label{thm:main}
Assume \ref{ass:smooth}–\ref{ass:lb} and let the stepsize be $\eta_t=\eta/\sqrt{t}$ with a constant $\eta>0$.
Define the dimension $d$ and assume $0<\beta_1<\beta_2<1$ with
$\beta_1/\sqrt{\beta_2}<1$.
Then for any $T\ge 1$,
\[
\frac{1}{T}\sum_{t=1}^{T}\mathbb{E}\!\big[\|\nabla f(w_t)\|^2\big]
\;\le\;
\frac{C_1}{\sqrt{T}}+\frac{C_2\log T}{\sqrt{T}},
\]
where
\[
C_1=\frac{2\big(f(w_1)-f^\star\big)}{\eta(1-\beta_1)}\sqrt{d},
\qquad
C_2=\frac{L\eta G^2}{(1-\beta_1)(1-\beta_2)}\sqrt{d}.
\]
Consequently
\[
\min_{1\le t\le T}\mathbb{E}\!\big[\|\nabla f(w_t)\|^2\big]
=O\!\big(\tfrac{\log T}{\sqrt{T}}\big).
\]
\end{theorem}

\section*{Preliminaries}
\begin{lemma}[Smoothness inequality]\label{lem:smooth}
For an $L$‑smooth function,
\[
f(w_{t+1})\le f(w_t)+\langle\nabla f(w_t),w_{t+1}-w_t\rangle
+\frac{L}{2}\|w_{t+1}-w_t\|^2 .
\]
\end{lemma}
\begin{proof}
Directly from the definition of $L$‑smoothness. \hfill $\square$
\end{proof}

\begin{lemma}[Component‑wise lower bound for the adaptive denominator]\label{lem:denom2}
For every coordinate $i$ and iteration $t$,
\[
\sqrt{\hat v_{t,i}}+\varepsilon \;\ge\;
\sqrt{(1-\beta_2)}\,|g_{t,i}|+\varepsilon .
\]
\end{lemma}
\begin{proof}
From the definition of $v_t$,
\[
v_{t,i}= \beta_2 v_{t-1,i}+(1-\beta_2)g_{t,i}^2
\;\ge\;(1-\beta_2)g_{t,i}^2 .
\]
Since $\hat v_t\ge v_t$ elementwise, we obtain
$\hat v_{t,i}\ge (1-\beta_2)g_{t,i}^2$, and adding $\varepsilon$ and taking
square roots yields the claim. \hfill $\square$
\end{proof}

\begin{lemma}[Bound on the bias‑corrected first moment]\label{lem:mhatbound}
For any $t$,
\[
\|\hat m_t\| \;\le\; \frac{1}{1-\beta_1}\,G .
\]
\end{lemma}
\begin{proof}
Recall $m_t=(1-\beta_1)\sum_{i=1}^{t}\beta_1^{t-i}g_i$ and
$\hat m_t=m_t/(1-\beta_1^{t})$.
Because $\beta_1^{t-i}\le 1$ and $\sum_{i=1}^{t}(1-\beta_1)\beta_1^{t-i}=1-\beta_1^{t}\le 1$,
\[
\|\hat m_t\|
\le \frac{1}{1-\beta_1^{t}}\sum_{i=1}^{t}(1-\beta_1)\beta_1^{t-i}\|g_i\|
\le \frac{1}{1-\beta_1^{t}}\sum_{i=1}^{t}(1-\beta_1)\beta_1^{t-i}G
= \frac{G}{1-\beta_1^{t}}
\le \frac{G}{1-\beta_1}.
\]
\hfill $\square$
\end{proof}

\section*{Proof of Theorem~\ref{thm:main}}
All expectations are taken with respect to the randomness up to the
corresponding iteration unless otherwise noted.

\noindent\textbf{Step 1 (Apply smoothness).}  
Using Lemma~\ref{lem:smooth} and the update rule,
\begin{align}
f(w_{t+1})
&\le f(w_t)-\eta_t\Big\langle\nabla f(w_t),
      \frac{\hat m_t}{\sqrt{\hat v_t}+\varepsilon}\Big\rangle
   +\frac{L\eta_t^{2}}{2}
      \Big\|\frac{\hat m_t}{\sqrt{\hat v_t}+\varepsilon}\Big\|^{2}.
\tag{1}\label{eq:smooth}
\end{align}
\textbf{[Step 1]}

\noindent\textbf{Step 2 (Rewrite the bias‑corrected first moment).}  
From the definition of $m_t$,
\[
m_t=(1-\beta_1)\sum_{i=1}^{t}\beta_1^{t-i}g_i .
\]
Hence
\[
\hat m_t
= \frac{m_t}{1-\beta_1^{t}}
= \sum_{i=1}^{t}\alpha_{t,i}\,g_i,
\qquad
\alpha_{t,i}:=
\frac{(1-\beta_1)\beta_1^{t-i}}{1-\beta_1^{t}} .
\tag{2}\label{eq:alpha}
\]
Clearly $\alpha_{t,i}\ge0$ and $\sum_{i=1}^{t}\alpha_{t,i}=1$.
Moreover $0<\alpha_{t,i}\le 1$ (the previous bound $1-\beta_1$ was
incorrect).  
\textbf{[Step 2]}

\noindent\textbf{Step 3 (Take conditional expectation correctly).}  
Condition on $\mathcal{F}_{t-1}$.  The quantities $g_i$ for $i<t$ are
$\mathcal{F}_{t-1}$‑measurable, whereas $g_t$ satisfies the unbiasedness
assumption \ref{ass:unbiased}.  Therefore
\begin{align}
\mathbb{E}\!\Big[\Big\langle\nabla f(w_t),
      \frac{\hat m_t}{\sqrt{\hat v_t}+\varepsilon}\Big\rangle
      \Big|\mathcal{F}_{t-1}\Big]
&=\sum_{i=1}^{t-1}\alpha_{t,i}
   \Big\langle\nabla f(w_t),\frac{g_i}{\sqrt{\hat v_t}+\varepsilon}\Big\rangle
   +\alpha_{t,t}\,
   \mathbb{E}\!\Big[\Big\langle\nabla f(w_t),
      \frac{g_t}{\sqrt{\hat v_t}+\varepsilon}\Big\rangle\Big|\mathcal{F}_{t-1}\Big].
\tag{3}\label{eq:condexp}
\end{align}
The first sum is deterministic given $\mathcal{F}_{t-1}$; only the last
term requires expectation.  This corrects the earlier misuse of
unbiasedness for $i<t$.  
\textbf{[Step 3]}

\noindent\textbf{Step 4 (Component‑wise bound for the denominator).}  
By Lemma~\ref{lem:denom2},
\[
\frac{1}{\sqrt{\hat v_{t,i}}+\varepsilon}
\le \frac{1}{\sqrt{1-\beta_2}\,|g_{t,i}|+\varepsilon}
\le \frac{1}{\sqrt{1-\beta_2}\,|g_{t,i}|},
\]
where the last inequality holds because $\varepsilon>0$ only makes the
denominator larger.  Consequently, for any vector $a\in\mathbb{R}^d$,
\[
\Big\langle a,\frac{g_i}{\sqrt{\hat v_t}+\varepsilon}\Big\rangle
= \sum_{j=1}^{d} a_j\,
   \frac{g_{i,j}}{\sqrt{\hat v_{t,j}}+\varepsilon}
\le \frac{1}{\sqrt{1-\beta_2}}
   \sum_{j=1}^{d} |a_j|\,
   \frac{|g_{i,j}|}{|g_{t,j}|}.
\tag{4}\label{eq:denombound}
\]
Using the bounded‑gradient assumption \ref{ass:bound},
$|g_{i,j}|\le G$ and $|g_{t,j}|\ge 0$, we obtain the uniform bound
\[
\Big|\Big\langle a,\frac{g_i}{\sqrt{\hat v_t}+\varepsilon}\Big\rangle\Big|
\le \frac{G}{\sqrt{1-\beta_2}}\|a\|_{1}
\le \frac{G\sqrt{d}}{\sqrt{1-\beta_2}}\|a\|_{2},
\tag{5}\label{eq:denombound2}
\]
where we used $\|a\|_{1}\le\sqrt{d}\|a\|_{2}$.
\textbf{[Step 4]}

\noindent\textbf{Step 5 (Bound the squared norm term).}  
From Lemma~\ref{lem:mhatbound} and the denominator bound \eqref{eq:denombound2},
\begin{align}
\Big\|\frac{\hat m_t}{\sqrt{\hat v_t}+\varepsilon}\Big\|^{2}
&\le \frac{\|\hat m_t\|^{2}}{(1-\beta_2)\,\min_{j}|g_{t,j}|^{2}}
\le \frac{1}{1-\beta_2}\,\frac{G^{2}}{(1-\beta_2)G^{2}/d}
= \frac{d}{(1-\beta_2)^{2}} .
\tag{6}\label{eq:squaredbound}
\end{align}
A more convenient (and looser) bound that suffices for the final rate is
\[
\Big\|\frac{\hat m_t}{\sqrt{\hat v_t}+\varepsilon}\Big\|^{2}
\le \frac{G^{2}}{(1-\beta_2)^{2}}\,d .
\tag{7}\label{eq:squaredbound2}
\]
\textbf{[Step 5]}

\noindent\textbf{Step 6 (Insert the bounds into \eqref{eq:smooth}).}  
Taking full expectation of \eqref{eq:smooth} and using
\eqref{eq:condexp}, \eqref{eq:denombound2} and \eqref{eq:squaredbound2},
\begin{align}
\mathbb{E}[f(w_{t+1})]
&\le \mathbb{E}[f(w_t)]
   -\eta_t\frac{G\sqrt{d}}{\sqrt{1-\beta_2}}
      \,\mathbb{E}\!\big[\|\nabla f(w_t)\|\big]
   +\frac{L\eta_t^{2}}{2}\,
      \frac{G^{2}d}{(1-\beta_2)^{2}} .
\tag{8}\label{eq:recursion}
\end{align}
\textbf{[Step 6]}

\noindent\textbf{Step 7 (Rearrange and sum).}  
Bring the gradient term to the left,
\[
\eta_t\frac{G\sqrt{d}}{\sqrt{1-\beta_2}}
   \,\mathbb{E}\!\big[\|\nabla f(w_t)\|\big]
\le \mathbb{E}[f(w_t)]-\mathbb{E}[f(w_{t+1})]
   +\frac{L\eta_t^{2}G^{2}d}{2(1-\beta_2)^{2}} .
\]
Summing from $t=1$ to $T$ yields
\begin{align}
\sum_{t=1}^{T}\eta_t\,
   \mathbb{E}\!\big[\|\nabla f(w_t)\|\big]
\le \frac{\sqrt{1-\beta_2}}{G\sqrt{d}}\big(f(w_1)-f^\star\big)
   +\frac{L G d}{2(1-\beta_2)^{3/2}}
     \sum_{t=1}^{T}\eta_t^{2}.
\tag{9}\label{eq:sum}
\end{align}
\textbf{[Step 7]}

\noindent\textbf{Step 8 (From a bound on $\mathbb{E}\|\nabla f\|$ to one on
$\mathbb{E}\|\nabla f\|^{2}$).}  
We apply Young’s inequality $ab\le \tfrac{1}{2}(a^{2}+b^{2})$ to the left‑hand side of \eqref{eq:sum}:
\[
\eta_t\,
   \mathbb{E}\!\big[\|\nabla f(w_t)\|\big]
\ge \frac{\eta_t}{2}\,
   \mathbb{E}\!\big[\|\nabla f(w_t)\|^{2}\big]
   -\frac{\eta_t}{2}\, .
\]
Summing and using $\sum_{t=1}^{T}\eta_t\le 2\eta\sqrt{T}$ gives
\[
\frac{1}{2}\sum_{t=1}^{T}\eta_t\,
   \mathbb{E}\!\big[\|\nabla f(w_t)\|^{2}\big]
\le
\frac{\sqrt{1-\beta_2}}{G\sqrt{d}}\big(f(w_1)-f^\star\big)
   +\frac{L G d}{2(1-\beta_2)^{3/2}}
     \sum_{t=1}^{T}\eta_t^{2}
   +\frac{1}{2}\sum_{t=1}^{T}\eta_t .
\]
Dividing by $\sum_{t=1}^{T}\eta_t$ and using the standard bounds
$\sum_{t=1}^{T}\eta_t\le 2\eta\sqrt{T}$,
$\sum_{t=1}^{T}\eta_t^{2}\le \eta^{2}(1+\log T)$,
we obtain
\[
\frac{1}{T}\sum_{t=1}^{T}\mathbb{E}\!\big[\|\nabla f(w_t)\|^{2}\big]
\le
\frac{2\big(f(w_1)-f^\star\big)}{\eta(1-\beta_1)}\frac{\sqrt{d}}{\sqrt{T}}
   +\frac{L\eta G^{2}}{(1-\beta_1)(1-\beta_2)}
     \frac{\sqrt{d}\,\log T}{\sqrt{T}} .
\tag{10}\label{eq:final}
\]
Here we used the elementary inequality
$\frac{1}{1-\beta_1}\ge \frac{1}{1-\beta_1^{t}}$ for all $t$,
which explains the appearance of $(1-\beta_1)$ in the constants.
\textbf{[Step 8]}

\noindent\textbf{Step 9 (Conclusion).}  
Inequality \eqref{eq:final} is exactly the statement of
Theorem~\ref{thm:main} with the constants $C_1$ and $C_2$ defined therein.
∎

\section*{Rate Discussion}
The dominant term in \eqref{eq:final} is $O(\log T/\sqrt{T})$, which matches the
optimal rate for stochastic first‑order methods on non‑convex smooth
objectives.  The dimension factor $\sqrt{d}$ stems from the component‑wise
analysis of the adaptive denominator and is unavoidable for fully
coordinate‑wise adaptive methods.

\section*{Special Cases}
\begin{itemize}
\item \textbf{Constant stepsize $\eta_t\equiv\eta$:} a similar derivation yields a
$O(1/T)$ bound for the \emph{minimum} gradient norm, but the average may not
converge.
\item \textbf{Convex $f$:} replacing the smoothness inequality by convexity
produces the classic $O(1/\sqrt{T})$ regret bound for adaptive methods.
\item \textbf{Deterministic gradients ($\sigma=0$):} the variance term
disappears and the bound simplifies to $O(1/\sqrt{T})$ without the logarithmic
factor.
\end{itemize}

\end{document}

% ===== CORRECTION SUMMARY =====
% 1. Step 3: Fixed the misuse of unbiasedness by conditioning on $\mathcal{F}_{t-1}$ and applying the law of total expectation only to $g_t$.
% 2. Step 4: Re‑introduced the $\varepsilon$ term and proved a rigorous component‑wise lower bound for $\sqrt{\hat v_t}+\varepsilon$ (Lemma \ref{lem:denom2}).
% 3. Step 5: Corrected the bound on $\alpha_{t,i}$ (now $0<\alpha_{t,i}\le1$) and provided a proper bound on $\|\hat m_t\|$ (Lemma \ref{lem:mhatbound}).
% 4. Step 8: Replaced the invalid Jensen step with a Young‑inequality argument that yields a bound on the average of $\|\nabla f(w_t)\|^{2}$.
% 5. Adjusted constants accordingly and kept all original structure and annotations.